% !TEX root = ../main.tex
\chapter{Proof of Inversion Lemmas} \label{app:inversion-lemmas}

Now we proof the inversion lemmas by simultaneous structural induction on derivation height.
\trulecases*
\begin{proof}
    The four cases are mutually exclusive by the unique outermost constructor of pre-type-expressions 
    (the BNF grammar for $E$ is unambiguous). We proceed by induction on $|\pi|$, with case analysis on 
    the last rule applied.

    \noindent
    \textbf{Group I: Direct Type Formation Rules} \\
    These immediately yield the required form and witnesses.
    \begin{itemize}
        \item[\ref{rule:prop}] We have that $A = \Prop$, this is Case 1.
        
        \item[\ref{rule:prop-type}] Then $A = \Proof(p)$ and $\Gamma \vdash p \colon \Prop$ as a sub-derivation. This is Case 2.
        \item[\ref{rule:prod-form}] Then $A = (\Pi x \colon B)C$ and $\Gamma, x \colon B \vdash C \type$ as a sub-derivation. This is Case 3.
        \item[\textsc{Introductory Rule}] Then, there exists a sort symbol $S \in \Sigma$ with an introductory rule $$z_1 \colon C_1, \dots, z_w \colon C_w \vdash S \type,$$ and derivable rules 
            $$\Gamma \vdash t_1 \colon C_1, \dots, \Gamma \vdash t_w \colon C_w[z_1 \leftarrow t_1, \dots, z_{w-1} \leftarrow t_{w-1}]$$ 
            such that $A = S[z_1 \leftarrow t_1, \dots, z_w \leftarrow t_w]$.
            This is Case 4.
    \end{itemize}

    \noindent
    \textbf{Group II: Context-Structural Rules} \\
    These change the context but not the outermost constructor of $A$; the case is transferred by applying the same rule to each sub-derivation.
    \begin{enumerate}
        \item[\ref{rule:equ}]
        By applying the induction hypothesis to $\Delta \vdash A \type$ (smaller $|\pi|$). One of the four cases holds in context $\Delta$. Transfer each to $\Gamma$:
        \begin{itemize}
            \item[Case 1]  Trivially holds for $\Gamma$.
            \item[Case 2]  Apply \ref{rule:equ} with $\Gamma = \Delta$ (from $\Delta = \Gamma$ by \ref{rule:csym}) to get $\Gamma \vdash p \colon \Prop$.
            \item[Case 3]  From $\Delta = \Gamma$ derive $\Delta, x \colon B = \Gamma, x \colon B$ by \ref{rule:cequ}. Apply \ref{rule:equ} to get $\Gamma, x \colon B \vdash C \type$.
            \item[Case 4]  Apply \ref{rule:equ} to each typing.
        \end{itemize}
        
        \item[\ref{rule:thin}]
        Apply the induction hypothesis to $\Gamma,\Delta \vdash A \type$. Transfer sub-derivations to $(\Gamma, y\colon B, \Delta)$ using \ref{rule:thin} (using \ref{rule:cont-thin} for the extended context in Case 3). All four cases transfer.
        
        \item[\ref{rule:sub}]
        Apply the induction hypothesis to $\Gamma, y \colon B, \Delta \vdash A \type$. Substitution preserves the outermost constructor.
        \begin{itemize}
            \item[Case 1]: $A[y \leftarrow s] = \Prop$.
            \item[Case 2]: $A[y\leftarrow s] = \Proof(p[y\leftarrow s])$. Apply \ref{rule:sub} to $\Gamma,y \colon B,\Delta \vdash p\colon\Prop$ to obtain $\Gamma,\Delta[y\leftarrow s] \vdash p[y\leftarrow s] \colon \Prop$.
            \item[Case 3]: $A[y\leftarrow s] = (\Pi x \colon C[y\leftarrow s])D[y\leftarrow s]$. Apply \ref{rule:sub} to $\Gamma,y \colon B,\Delta,x \colon C \vdash D \type$ to obtain $\Gamma,\Delta[y\leftarrow s],x \colon C[y\leftarrow s] \vdash D[y\leftarrow s] \type$.
            \item[Case 4]: $A[y\leftarrow s] = S[z_i \leftarrow t_i[y\leftarrow s]]$. Apply \ref{rule:sub} to each typing derivation using $(C_i[z_1\leftarrow t_1,\ldots])[y\leftarrow s] = C_i[z_1 \leftarrow t_1[y\leftarrow s],\ldots]$.
        \end{itemize}
    \end{enumerate}

    \noindent
    \textbf{Group III: The Reflection and Equality Rules}
    \begin{itemize}
        \item[\ref{rule:reflection2}] 
        The sub-derivation $\pi' \colon \Gamma \vdash t \colon A$ satisfies $|\pi'| < |\pi|$. Apply Lemma \ref{lemma:in-rule-cases} by induction hypothesis to $\pi'$, we get 5 cases:
        \begin{itemize}
            \item[Case 1] Let $\Gamma = \Gamma_0, x\colon A, \Gamma_1$. Context $\Gamma$ was formed by \ref{rule:cont-intro} from $\Gamma_0 \vdash A \type$. By applying \ref{rule:thin}, $\Gamma \vdash A \type$ is derivable by a shorter derivation. Apply the IH.
            \item[Case 2] From $\Gamma,x\colon B \vdash s\colon C$, apply \ref{rule:reflection2} (strictly smaller height) to get $\Gamma,x\colon B \vdash C \type$. This is Case 3.
            \item[Case 3] From $\Gamma,x\colon B \vdash C \type$ and $\Gamma \vdash r\colon B$, apply \ref{rule:sub} to obtain $\Gamma \vdash C[x\leftarrow r] \type$ by a shorter derivation. Apply the IH.
            \item[Case 4] This is Case 1.
            \item[Case 5] Apply \ref{rule:reflection2} to $\Gamma \vdash F[z_i\leftarrow t_i] \colon D[z_i\leftarrow t_i]$ (shorter) to get $\Gamma \vdash A \type$. Apply the IH.
        \end{itemize}

        \item [\ref{rule:refl-typ}]
        $A$ is syntactically of forms (1)-(4). We proceed by inner induction on the derivation of $\Gamma \vdash A = A$:
        \begin{itemize}
            \item[\ref{rule:refl-typ}] The derivation of $\Gamma \vdash A \type$ is strictly shorter; apply the outer induction hypothesis.
            \item[\ref{rule:prod-equ}] From $\Gamma,x\colon B \vdash C = C$, apply \ref{rule:refl-typ} (IH) to extract $\Gamma,x\colon B \vdash C \type$. This is Case 3.
            \item[\ref{rule:prop-equ}] From $\Gamma \vdash p = p \colon \Prop$, apply \ref{rule:refl-obj} to extract $\Gamma \vdash p \colon \Prop$. This is Case 2.
            \item[\ref{rule:forall-elim}] $\Gamma \vdash \Proof((\forall x\colon B)p) = (\Pi x\colon B)\Proof(p)$. If $A = \Proof((\forall x\colon B)p)$, it is Case 2. If $A \equiv (\Pi x\colon B)\Proof(p)$, it is Case 3. \qedhere
        \end{itemize}
    \end{itemize}
\end{proof}

\inrulecases*
\begin{proof}
    By simultaneous induction on $|\pi|$ with Lemma \ref{lemma:t-rule-cases}.

    \noindent
    \textbf{Group I: Direct Term Introduction Rules}
    \begin{itemize}
        \item[\ref{rule:ass}] $t = x \in \Var(\Gamma)$. This is Case 1.
        
        \item[\ref{rule:prod-intro}] $t = (\lambda x\colon B)s$ and $A = (\Pi x\colon B)C$. The premise is a strict sub-derivation. This is Case 2.
        
        \item[\ref{rule:prod-elim}] $t = \App([x\colon B]C,f,r)$ and $A = C[x\leftarrow r]$. Premises are strict sub-derivations. This is Case 3.
        
        \item[\ref{rule:forall-intro}] $t = (\forall x\colon B)p$ and $A = \Prop$. This is Case 4.
        
        \item[\textsc{Introductory Rule}] For an operator $F \in \Sigma$,
        $t = F[z_1 \leftarrow t_1, \dots, z_w \leftarrow t_w]$ with all premises as strict sub-derivations. This is Case 5.
    \end{itemize}

    \noindent
    \textbf{Group II: Context-Structural Rules} \\
    These preserve the syntactic form of $t$; only the context and type representative change.
    \begin{itemize}
        \item[\ref{rule:equ}] Apply the induction hypothesis to $\Delta \vdash t \colon A$. Transfer each sub-derivation from context $\Delta$ to $\Gamma$ using \ref{rule:equ} (and \ref{rule:cequ} where necessary). All cases transfer safely.
        
        \item[\ref{rule:thin}] Apply the induction hypothesis to $\Gamma,\Delta \vdash t\colon A$. Transfer sub-derivations using \ref{rule:thin} (using \ref{rule:cont-thin} for extended contexts), valid by \ref{rule:prop}. All cases transfer.
        
        \item[\ref{rule:sub}] Apply the induction hypothesis to $\Gamma,y\colon B,\Delta \vdash t\colon A$. Substitution is transparent to the outermost constructor of $t$. If $t = x = y$, then $t[y\leftarrow s] = s$. Extend $\Gamma \vdash s\colon B$ by \ref{rule:thin} and apply the induction hypothesis to $\Gamma \vdash s\colon B$ (smaller height) to recursively determine the form of $s$.
    \end{itemize}

    \noindent
    \textbf{Group III: Type-Conversion and Equality Rules}
    \begin{itemize}
        \item[\ref{rule:conv1}]
        Apply the induction hypothesis to $\Gamma \vdash t \colon A'$ (smaller height). The syntactic form of $t$ is unchanged; we propagate the type claim from $A'$ to $A$ using provable equality $\Gamma \vdash A' = A$ and \ref{rule:trans-typ}. For example, in Case 2, $A'$ is provably equal to $(\Pi x\colon B)C$, and by transitivity, $A$ is provably equal to $(\Pi x\colon B)C$. 
        
        \item[\ref{rule:refl-obj}] 
        By inner induction on the derivation of term equality. If derived from $\Gamma \vdash t\colon A$, it is shorter (apply outer induction hypothesis). If derived by \ref{rule:trans-obj}, the companion term-equality lemma determines the form. If derived by \ref{rule:beta-rule} or \ref{rule:eta-rule}, $t$ maps directly to Case 3 or Case 2, respectively. \qedhere
    \end{itemize}
\end{proof}